\part{VII--VIII klasė}

\section{Racionalieji skaičiai ir laipsniai}\label{sec:racionalieji-skaiciai-ir-laipsniai}

\subsection{Apibrėžtys}
\begin{apibreztis}
Racionalusis skaičius yra skaičius, kuri galima užrašyti trupmena $\frac{a}{b}$, kai $a,b\in\mathbb{Z}$ ir $b\ne0$. Racionalieji skaičiai sudaro aibė $\mathbb{Q}$. Pavyzdziui, $-7=\frac{-7}{1}$, $0{,}25=\frac14$, o $3{,}(6)=\frac{11}{3}$.
\end{apibreztis}
\begin{apibreztis}
Laipsnis $a^n$, kai $n\in\mathbb{N}$, reiškia $n$ vienodų dauginamųjų sandauga:
\[
a^n=\underbrace{a\cdot a\cdot\ldots\cdot a}_{n\ \text{kartu}}.
\]
Skaicius $a$ vadinamas pagrindu, $n$ -- laipsnio rodikliu.
\end{apibreztis}
\begin{apibreztis}
Kai $a\ne0$, apibrėžiama $a^0=1$, o neigiamas rodiklis reiškia atvirkstini skaičių:
\[
a^{-n}=\frac{1}{a^n},\qquad n\in\mathbb{N}.
\]
\end{apibreztis}
\begin{apibreztis}
Standartine skaičiaus išraiška yra pavidalas $a\cdot10^n$, kai $1\le |a|<10$, o $n$ yra sveikasis skaičius. Tokiu pavidalu patogu rasyti labai didelius ir labai mažus skaičius.
\end{apibreztis}

\subsection{Teoremos ir savybės}
\begin{savybe}
Kai $a\ne0$, $b\ne0$, o $m,n\in\mathbb{Z}$, galioja laipsniu savybes:
\[
a^m\cdot a^n=a^{m+n},\qquad \frac{a^m}{a^n}=a^{m-n},\qquad (a^m)^n=a^{mn}.
\]
\end{savybe}
\begin{savybe}
Vienodu rodikliu laipsniai dauginami ir dalijami taip:
\[
a^n b^n=(ab)^n,\qquad \frac{a^n}{b^n}=\left(\frac{a}{b}\right)^n.
\]
\end{savybe}
\begin{savybe}
Jei $n$ lyginis, tai $(-a)^n=a^n$; jei $n$ nelyginis, tai $(-a)^n=-a^n$. Be skliaustų užrašas $-a^n$ reiškia $-(a^n)$.
\end{savybe}
\begin{pastaba}
Veiksmu tvarka: skliaustai, laipsniai, daugyba ir dalyba iš kairės i dešinę, sudėtis ir atimtis iš kairės i dešinę.
\end{pastaba}

\subsection{Taisyklės ir algoritmai}
\begin{enumerate}[label=\arabic*.]
\item Sudedant arba atimant racionaliuosius skaičius, trupmenas suveskite i bendra vardikli, atlikite veiksma su skaitikliais ir suprastinkite atsakymą.
\item Dauginant trupmenas, dauginkite skaitiklius ir vardiklius; dalijant iš trupmenos, dauginkite iš atvirkstines trupmenos.
\item Pertvarkydami laipsnius, pirmiausia sugrupuokite vienodus pagrindus arba vienodus rodiklius, tada taikykite savybes.
\item Skaiciu rasydami standartine išraiška, kableli perkelkite taip, kad pries ji liktu vienas nenulinis skaitmuo; perkeltu vietu skaičius tampa $10$ rodikliu.
\end{enumerate}
\paragraph{Dažnos klaidos.}
Nereikia painioti $(-3)^2=9$ ir $-3^2=-9$. Taip pat $a^{-2}$ nėra neigiamas skaičius; tai $\frac{1}{a^2}$, jei $a\ne0$. Dar viena klaida -- sudedant trupmenas sudeti vardiklius: iš tikruju $\frac12+\frac13=\frac56$, ne $\frac25$.

\subsection{Iliustruojantys pavyzdžiai}
\textbf{1 pavyzdys.} Apskaiciuokime
\[
-\frac34+\frac56.
\]
Bendras vardiklis yra $12$, todėl
\[
-\frac34+\frac56=-\frac9{12}+\frac{10}{12}=\frac1{12}.
\]
\textbf{Atsakymas: $\frac1{12}$.}

\textbf{2 pavyzdys.} Supaprastinkime
\[
\frac{2^5\cdot2^{-2}}{2^3}.
\]
Vienodas pagrindas $2$, todėl rodiklius sudedame arba atimame:
\[
\frac{2^5\cdot2^{-2}}{2^3}=2^{5+(-2)-3}=2^0=1.
\]
\textbf{Atsakymas: $1$.}

\textbf{3 pavyzdys.} Skaiciu $0{,}00072$ užrašykime standartine išraiška.
Kableli perkeliame $4$ vietomis i dešinę, kad gautume $7{,}2$. Kad skaičiaus verte nepasikeistu, dauginame iš $10^{-4}$:
\[
0{,}00072=7{,}2\cdot10^{-4}.
\]
\textbf{Atsakymas: $7{,}2\cdot10^{-4}$.}

\textbf{4 pavyzdys.} Palyginkime $4{,}1\cdot10^6$ ir $3{,}9\cdot10^7$.
Suvienodiname rodikli:
\[
3{,}9\cdot10^7=39\cdot10^6.
\]
Kadangi $39\cdot10^6>4{,}1\cdot10^6$, antras skaičius didesnis.
\textbf{Atsakymas: $3{,}9\cdot10^7$ didesnis.}

\paragraph{Pasitikrinkite.}
\begin{enumerate}[label=\alph*)]
\item Apskaiciuokite $\frac25-\frac76$.
\item Supaprastinkite $(3^2)^3\cdot3^{-4}$.
\item Uzrasykite $58000000$ standartine išraiška.
\end{enumerate}

\subsection{Ryšys su kitomis temomis}
Racionalieji skaiciai tęsia trupmenu ir sveikuju skaiciu temas iš \ref{sec:sveikieji-skaiciai-ir-racionalieji-skaiciai}. Laipsniai reikalingi saknims \ref{sec:saknys-ir-realiojo-skaiciaus-samprata}, algebriniams pertvarkiams \ref{sec:algebriniu-reiskiniu-pertvarkiai-ir-skaidymas}, standartinei israiskai fizikoje ir velesniems realiuju skaiciu intervalams \ref{sec:realieji-skaiciai-ir-intervalai}.

\section{Šaknys ir realiojo skaičiaus samprata}\label{sec:saknys-ir-realiojo-skaiciaus-samprata}

\subsection{Apibrėžtys}
\begin{apibreztis}
Kvadratine šaknis iš neneigiamo skaičiaus $a$ yra toks neneigiamas skaičius $\sqrt a$, kurio kvadratas lygus $a$:
\[
(\sqrt a)^2=a,\qquad a\ge0.
\]
\end{apibreztis}
\begin{apibreztis}
Kubine šaknis iš skaičiaus $a$ yra toks skaičius $\sqrt[3]{a}$, kurio kubas lygus $a$. Kubine šaknis apibrėžiama ir neigiamiems skaičiams, pavyzdžiui, $\sqrt[3]{-8}=-2$.
\end{apibreztis}
\begin{apibreztis}
Iracionalusis skaičius yra skaičius, kurio negalima užrašyti trupmena $\frac{a}{b}$, kai $a,b\in\mathbb{Z}$, $b\ne0$. Realieji skaičiai yra racionalieji ir iracionalieji skaičiai kartu; ju aibė žymima $\mathbb{R}$.
\end{apibreztis}
\begin{apibreztis}
Skaiciu aibė yra skaičių rinkinys. Jei kiekvienas aibės $A$ elementas priklauso aibei $B$, tai $A$ yra $B$ poaibis ir rasome $A\subset B$.
\end{apibreztis}

\subsection{Teoremos ir savybės}
\begin{savybe}
Jei $a\ge0$ ir $b\ge0$, tai
\[
\sqrt{ab}=\sqrt a\cdot\sqrt b.
\]
Jei $a\ge0$ ir $b>0$, tai
\[
\sqrt{\frac{a}{b}}=\frac{\sqrt a}{\sqrt b}.
\]
\end{savybe}
\begin{savybe}
Visada $\sqrt{a^2}=|a|$. Lygybe $\sqrt{a^2}=a$ teisinga tik tada, kai $a\ge0$.
\end{savybe}
\begin{savybe}
Kubinei šakniai galioja $\sqrt[3]{ab}=\sqrt[3]{a}\sqrt[3]{b}$ ir $\sqrt[3]{a^3}=a$ visiems realiesiems $a$.
\end{savybe}
\begin{pastaba}
Skaiciu aibės susijusios taip:
\[
\mathbb{N}\subset\mathbb{Z}\subset\mathbb{Q}\subset\mathbb{R}.
\]
Iracionalieji skaičiai priklauso $\mathbb{R}$, bet nepriklauso $\mathbb{Q}$.
\end{pastaba}

\subsection{Taisyklės ir algoritmai}
\begin{enumerate}[label=\arabic*.]
\item Norint supaprastinti kvadratinę šaknį, pošakninį skaičių išskaidykite taip, kad vienas daugiklis butu didziausias galimas kvadratas.
\item Norint įvertinti $\sqrt a$, raskite gretimus sveikuosius skaičius, kuriu kvadratai riboja $a$.
\item I po šaknies ženklą teigiama daugikli galima ikelti kvadratu: $k\sqrt a=\sqrt{k^2a}$, kai $k\ge0$.
\item Priskirdami skaičių aibei, pirmiausia klauskite, ar jis natūralusis, tada sveikasis, racionalusis, realusis.
\end{enumerate}
\paragraph{Dažnos klaidos.}
Negalima rasyti $\sqrt{a+b}=\sqrt a+\sqrt b$: pavyzdžiui, $\sqrt{9+16}=5$, bet $\sqrt9+\sqrt{16}=7$. Taip pat $\sqrt{49}$ yra $7$, ne $\pm7$; lygtis $x^2=49$ turi du sprendinius $x=\pm7$.

\subsection{Iliustruojantys pavyzdžiai}
\textbf{1 pavyzdys.} Supaprastinkime $\sqrt{72}$.
\[
72=36\cdot2,\qquad \sqrt{72}=\sqrt{36}\sqrt2=6\sqrt2.
\]
\textbf{Atsakymas: $6\sqrt2$.}

\textbf{2 pavyzdys.} Tarp kokiu sveikųjų skaičių yra $\sqrt{50}$?
\[
7^2=49,\qquad 8^2=64,\qquad 49<50<64.
\]
Vadinasi,
\[
7<\sqrt{50}<8.
\]
\textbf{Atsakymas: tarp $7$ ir $8$.}

\textbf{3 pavyzdys.} Apskaiciuokime $\sqrt{(-11)^2}$.
\[
\sqrt{(-11)^2}=|-11|=11.
\]
\textbf{Atsakymas: $11$.}

\textbf{4 pavyzdys.} Nustatykime, kurioms aibems priklauso skaičiai $-5$, $\frac23$, $\sqrt9$, $\sqrt2$.
Turime $\sqrt9=3$, todėl jis priklauso $\mathbb{N}$, $\mathbb{Z}$, $\mathbb{Q}$ ir $\mathbb{R}$. Skaicius $-5$ priklauso $\mathbb{Z}$, $\mathbb{Q}$, $\mathbb{R}$. Skaicius $\frac23$ priklauso $\mathbb{Q}$ ir $\mathbb{R}$. Skaicius $\sqrt2$ yra iracionalusis, todėl priklauso tik $\mathbb{R}$ iš siu aibiu.

\paragraph{Pasitikrinkite.}
\begin{enumerate}[label=\alph*)]
\item Supaprastinkite $\sqrt{98}$.
\item Ivertinkite, tarp kokiu sveikųjų skaičių yra $\sqrt[3]{30}$.
\item Paaiskinkite, kodel $\sqrt{0{,}25}$ yra racionalusis skaičius.
\end{enumerate}

\subsection{Ryšys su kitomis temomis}
Saknys remiasi laipsniais iš \ref{sec:racionalieji-skaiciai-ir-laipsniai}. Jos naudojamos Pitagoro teoremoje \ref{sec:pitos-teorema}, erdviniu kunu skaiciavimuose \ref{sec:erdviniai-kunai-ir-ju-pjuviai} ir velesneje realiuju skaiciu bei funkciju teorijoje \ref{sec:realieji-skaiciai-ir-intervalai}.

\section{Procentai, proporcijos ir mastelis}\label{sec:procentai-proporcijos-ir-mastelis}

\subsection{Apibrėžtys}
\begin{apibreztis}
Procentas yra viena simtoji dalis: $1\%=\frac1{100}=0{,}01$.
\end{apibreztis}
\begin{apibreztis}
Proporcija yra dvieju santykiu lygybe
\[
\frac{a}{b}=\frac{c}{d},\qquad b\ne0,\ d\ne0.
\]
\end{apibreztis}
\begin{apibreztis}
Mastelis yra brėžinyje, plane ar žemėlapyje pavaizduoto ilgio ir tikrojo ilgio santykis. Mastelis $1:500$ reiškia, kad $1$ brėzinio vienetas atitinka $500$ tokiu paciu tikruju vienetu.
\end{apibreztis}
\begin{apibreztis}
Procentinis pokytis rodo, keliais procentais nauja reikšmė skiriasi nuo pradinės:
\[
\text{pokytis procentais}=\frac{\text{nauja reikšmė}-\text{pradine reikšmė}}{\text{pradine reikšmė}}\cdot100\%.
\]
\end{apibreztis}

\subsection{Teoremos ir savybės}
\begin{savybe}
Proporcijoje $\frac{a}{b}=\frac{c}{d}$ kryzmines sandaugos lygios:
\[
ad=bc.
\]
\end{savybe}
\begin{savybe}
Padidinus dydį $p\%$, jis dauginamas iš $1+\frac{p}{100}$; sumazinus $p\%$, dauginamas iš $1-\frac{p}{100}$.
\end{savybe}
\begin{savybe}
Jei dydis kelis kartus iš eiles keiciamas procentais, bendras daugiklis gaunamas sudauginus atskirus pokycio daugiklius.
\end{savybe}
\begin{pastaba}
Du pokyciai $+20\%$ ir $-20\%$ vienas kito nepanaikina: $1{,}2\cdot0{,}8=0{,}96$, vadinasi, lieka $4\%$ sumazejimas.
\end{pastaba}

\subsection{Taisyklės ir algoritmai}
\begin{enumerate}[label=\arabic*.]
\item $p\%$ skaičiaus $N$ randama skaičiuojant $\frac{p}{100}\cdot N$.
\item Jei zinoma dalis ir procentas, visa reikšmė randama dalijant dali iš procento desimtaines išraiškos.
\item Proporcijoje nezinoma nari raskite kryzminiu sandaugu lygybe.
\item Mastelio uzdaviniuose pirma suvienodinkite matavimo vienetus, tada naudokite santyki $\frac{\text{brezinio ilgis}}{\text{tikrasis ilgis}}$.
\end{enumerate}
\paragraph{Dažnos klaidos.}
Procentus reikia skaičiuoti nuo tinkamos bazės. Jei prekė pirma atpigo, antras pokytis skaičiuojamas nuo naujos kainos. Mastelyje dazna klaida -- maisyti centimetrus ir metrus; santykyje abu ilgiai turi būti tais paciais vienetais.

\subsection{Iliustruojantys pavyzdžiai}
\textbf{1 pavyzdys.} Preke kainavo $80$ euru. Kaina sumazinta $15\%$. Raskime nauja kaina.
\[
80\cdot\left(1-\frac{15}{100}\right)=80\cdot0{,}85=68.
\]
\textbf{Atsakymas: $68$ eurai.}

\textbf{2 pavyzdys.} Ispręskime proporcija
\[
\frac{x}{12}=\frac58.
\]
Pagal kryzmines sandaugas:
\[
8x=12\cdot5=60,\qquad x=7{,}5.
\]
\textbf{Atsakymas: $x=7{,}5$.}

\textbf{3 pavyzdys.} Zemelapio mastelis $1:20000$. Atstumas žemėlapyje $4$ cm. Raskime tikraji atstuma.
\[
4\cdot20000=80000\ \text{cm}=800\ \text{m}=0{,}8\ \text{km}.
\]
\textbf{Atsakymas: $0{,}8$ km.}

\textbf{4 pavyzdys.} Kaina padidinta $10\%$, o po to sumazinta $10\%$. Pradine kaina $50$ euru. Kokia galutine kaina?
\[
50\cdot1{,}10\cdot0{,}90=49{,}5.
\]
\textbf{Atsakymas: $49{,}50$ euro.}

\paragraph{Pasitikrinkite.}
\begin{enumerate}[label=\alph*)]
\item Raskite $18\%$ skaičiaus $250$.
\item Skaicius padidejo nuo $40$ iki $46$. Keliais procentais jis padidejo?
\item Plane $3$ cm atitinka $12$ m. Koks mastelis?
\end{enumerate}

\subsection{Ryšys su kitomis temomis}
Procentai ir proporcijos remiasi trupmenomis \ref{sec:paprastosios-ir-desimtaines-trupmenos} ir tiesioginiu proporcingumu. Jie tiesiogiai naudojami finansiniuose skaiciavimuose \ref{sec:finansiniai-skaiciavimai-palukanos-kreditas-lizingas}, panasiuose trikampiuose \ref{sec:panasus-trikampiai-ir-talio-teorema}, mastelyje \ref{sec:mastelis-breziniai-ir-transformacijos} ir statistikoje \ref{sec:statistika-imtys-ir-sklaida}.

\section{Algebriniai reiškiniai ir daugianariai}\label{sec:algebriniai-reiskiniai-ir-daugianariai}

\subsection{Apibrėžtys}
\begin{apibreztis}
Algebrinis reiškinys yra skaičių, raidziu ir veiksmu zenklu junginys, turintis reikšmė, kai raidemis pazymetiems kintamiesiems parenkamos leistinos reikšmes.
\end{apibreztis}
\begin{apibreztis}
Vienanaris yra skaičiaus ir raidziu laipsniu sandauga, pavyzdžiui, $-5a^2b$. Vienanario skaicine dalis vadinama koeficientu.
\end{apibreztis}
\begin{apibreztis}
Daugianaris yra vienanariu suma. Dvinaris turi du narius, trinaris -- tris narius. Daugianario nariai yra panasūs, jei ju raidine dalis su tais paciais rodikliais sutampa.
\end{apibreztis}
\begin{apibreztis}
Reiškinio reikšmių sritis šiame koncentre paprastai reiškia kintamųjų reikšmes, kurioms reiškinys turi prasmę, pavyzdžiui, trupmenoje vardiklis negali būti lygus nuliui.
\end{apibreztis}

\subsection{Teoremos ir savybės}
\begin{savybe}
Daugyba skirstoma sudeciai ir atimciai:
\[
a(b+c)=ab+ac,\qquad a(b-c)=ab-ac.
\]
\end{savybe}
\begin{savybe}
Panasieji nariai sudedami sudedant ju koeficientus, o raidine dalis paliekama ta pati:
\[
3x^2-7x^2=-4x^2.
\]
\end{savybe}
\begin{savybe}
Dažnosios daugybos formulės:
\[
(a+b)^2=a^2+2ab+b^2,
\]
\[
(a-b)^2=a^2-2ab+b^2,\qquad (a-b)(a+b)=a^2-b^2.
\]
\end{savybe}

\subsection{Taisyklės ir algoritmai}
\begin{enumerate}[label=\arabic*.]
\item Reiškinio reiksmei rasti vietoje raidziu įrašykite skaičius, islaikykite skliaustus ir veiksmu tvarka.
\item Sutraukdami daugianari, pirmiausia atskliauskite, tada sugrupuokite panašiuosius narius.
\item Dauginant du dvinarius, kiekviena pirmojo skliausto nari dauginkite iš kiekvieno antrojo skliausto nario.
\item Pries naudodami formule, patikrinkite jos sandara: $(a-b)^2$ turi vidurini nari $-2ab$, o ne $-ab$.
\end{enumerate}
\paragraph{Dažnos klaidos.}
Klaidinga manyti, kad $(a+b)^2=a^2+b^2$. Truksta vidurinio nario $2ab$. Taip pat negalima sudeti nepanasiu nariu: $2x+3y$ nėra $5xy$.

\subsection{Iliustruojantys pavyzdžiai}
\textbf{1 pavyzdys.} Sutraukime reiškinį
\[
3x-2(4-x)+5.
\]
Atskliaudziame ir sutraukiame:
\[
3x-8+2x+5=5x-3.
\]
\textbf{Atsakymas: $5x-3$.}

\textbf{2 pavyzdys.} Sudauginkime
\[
(2a-3)(a+5).
\]
\[
(2a-3)(a+5)=2a^2+10a-3a-15=2a^2+7a-15.
\]
\textbf{Atsakymas: $2a^2+7a-15$.}

\textbf{3 pavyzdys.} Raskime reiškinio $x^2-4x+7$ reikšmė, kai $x=-2$.
\[
(-2)^2-4(-2)+7=4+8+7=19.
\]
\textbf{Atsakymas: $19$.}

\paragraph{Pasitikrinkite.}
\begin{enumerate}[label=\alph*)]
\item Sutraukite $5a-3(2a-4)+7$.
\item Sudauginkite $(x+4)(x-6)$.
\item Isskleiskite $(2m-3)^2$.
\end{enumerate}

\subsection{Ryšys su kitomis temomis}
Algebriniai reiskiniai tęsia raidiniu reiskiniu temą iš \ref{sec:skaitiniai-ir-raidiniai-reiskiniai}. Jie reikalingi lygtims \ref{sec:tiesines-lygtys-nelygybes-ir-sistemos}, funkciju formulems \ref{sec:tiesine-funkcija-ir-grafikas}, daugianariu skaidymui \ref{sec:algebriniu-reiskiniu-pertvarkiai-ir-skaidymas} ir kvadratinems lygtims \ref{sec:kvadratines-lygtys-ir-nelygybes}.

\section{Tiesinės lygtys, nelygybės ir sistemos}\label{sec:tiesines-lygtys-nelygybes-ir-sistemos}

\subsection{Apibrėžtys}
\begin{apibreztis}
Tiesine lygtis su vienu nežinomuoju yra lygtis, kuria galima pertvarkyti i pavidala $ax+b=0$, kur $a\ne0$.
\end{apibreztis}
\begin{apibreztis}
Tiesine nelygybė su vienu nežinomuoju yra nelygybė, kuria galima pertvarkyti i viena iš pavidalu $ax+b<0$, $ax+b>0$, $ax+b\le0$, $ax+b\ge0$.
\end{apibreztis}
\begin{apibreztis}
Lygtis su dviem nežinomaisiais turi pavidalo $ax+by=c$ sprendinius -- skaičių poras $(x;y)$, kurios ja paverčia teisinga lygybe.
\end{apibreztis}
\begin{apibreztis}
Tiesiniu lygciu sistema yra kelios lygtys, kurias turi tenkinti ta pati nezinomuju pora.
\end{apibreztis}

\subsection{Teoremos ir savybės}
\begin{savybe}
Lygties abi puses galima prideti, atimti, dauginti arba dalyti iš to paties nelygaus nuliui skaičiaus; sprendiniu aibė nepasikeicia.
\end{savybe}
\begin{savybe}
Dauginant arba dalijant nelygybės abi puses iš neigiamo skaičiaus, nelygybės ženklas keiciamas priesingu.
\end{savybe}
\begin{savybe}
Tiesiniu lygciu sistemos su dviem nežinomaisiais sprendinys grafiskai yra dvieju tiesiu sankirtos taskas. Sistema gali tureti viena sprendini, be galo daug sprendiniu arba netureti sprendiniu.
\end{savybe}

\subsection{Taisyklės ir algoritmai}
\begin{enumerate}[label=\arabic*.]
\item Sprendziant tiesinė lygtį, atskliauskite, sutraukite panašiuosius narius, nežinomojo narius perkelkite i viena puse, skaičius -- i kita, tada padalykite iš koeficiento.
\item Sprendziant nelygybė, atlikite tuos pacius pertvarkius kaip lygtyje, bet stebekite daugyba ar dalyba iš neigiamo skaičiaus.
\item Sistema keitimo budu: iš vienos lygties išreikškite viena nezinomaji, įrašykite i kita lygtį, raskite pirma nezinomaji, tada grizkite prie antro.
\item Sistema sudeties budu: lygtys dauginamos taip, kad sudejus isnyktu vienas nezinomasis.
\end{enumerate}
\paragraph{Dažnos klaidos.}
Perkeliant nari i kita lygties puse, keiciamas jo ženklas. Nelygybese dažniausia klaida -- pamiršti pakeisti ženklą dalijant iš neigiamo skaičiaus. Sistemose reikia pateikti pora $(x;y)$, o ne du nesusietus skaičius.

\subsection{Iliustruojantys pavyzdžiai}
\textbf{1 pavyzdys.} Ispręskime lygtį
\[
3(2x-1)=4x+9.
\]
\[
6x-3=4x+9,\qquad 2x=12,\qquad x=6.
\]
\textbf{Atsakymas: $x=6$.}

\textbf{2 pavyzdys.} Ispręskime nelygybė
\[
-2x+5>11.
\]
\[
-2x>6.
\]
Dalyba iš $-2$ keicia ženklą:
\[
x<-3.
\]
\textbf{Atsakymas: $x<-3$.}

\textbf{3 pavyzdys.} Ispręskime sistema
\[
\begin{cases}
x+y=9,\\
x-y=3.
\end{cases}
\]
Sudedame lygtis:
\[
2x=12,\qquad x=6.
\]
Tada $6+y=9$, todėl $y=3$.
\textbf{Atsakymas: $(6;3)$.}

\textbf{4 pavyzdys.} Sistema
\[
\begin{cases}
2x+y=7,\\
y=3x-8
\end{cases}
\]
sprendžiama keitimo budu:
\[
2x+(3x-8)=7,\qquad 5x=15,\qquad x=3.
\]
Tada $y=3\cdot3-8=1$.
\textbf{Atsakymas: $(3;1)$.}

\paragraph{Pasitikrinkite.}
\begin{enumerate}[label=\alph*)]
\item Ispręskite $5(x-2)=2x+8$.
\item Ispręskite $3-4x\le 15$.
\item Ispręskite sistema $\begin{cases}x+2y=11,\\ x-y=2.\end{cases}$
\end{enumerate}

\subsection{Ryšys su kitomis temomis}
Tiesines lygtys remiasi paprastu lygciu sprendimu \ref{sec:lygtys-ir-nelygybes-su-vienu-nezinomuoju}. Sistemos siejasi su tiesines funkcijos grafikais \ref{sec:tiesine-funkcija-ir-grafikas}, nelygybiu sprendiniu aibemis \ref{sec:tiesines-nelygybes-ir-sprendiniu-aibes} ir realaus turinio modeliavimu.

\section{Tiesinė funkcija ir grafikas}\label{sec:tiesine-funkcija-ir-grafikas}

\subsection{Apibrėžtys}
\begin{apibreztis}
Funkcija yra taisykle, pagal kuria kiekvienai leistinai $x$ reiksmei priskiriama viena $y$ reikšmė.
\end{apibreztis}
\begin{apibreztis}
Tiesine funkcija yra funkcija
\[
y=kx+b,
\]
kur $k$ ir $b$ yra skaičiai. Skaicius $k$ vadinamas krypties koeficientu, o $b$ -- pradine reikšmė arba susikirtimo su $y$ asimi ordinate.
\end{apibreztis}
\begin{apibreztis}
Funkcijos nulis yra tokia $x$ reikšmė, su kuria $y=0$.
\end{apibreztis}
\begin{apibreztis}
Tiesinis sarysis yra priklausomybe, kurios pokytis vienam $x$ vienetui yra pastovus.
\end{apibreztis}

\subsection{Teoremos ir savybės}
\begin{savybe}
Funkcijos $y=kx+b$ grafikas yra tiese.
\end{savybe}
\begin{savybe}
Jei $k>0$, funkcija dideja; jei $k<0$, funkcija mazeja; jei $k=0$, funkcija yra pastovi.
\end{savybe}
\begin{savybe}
Tieses krypties koeficientas per du taškus $(x_1;y_1)$ ir $(x_2;y_2)$, kai $x_1\ne x_2$, yra
\[
k=\frac{y_2-y_1}{x_2-x_1}.
\]
\end{savybe}

\subsection{Taisyklės ir algoritmai}
\begin{enumerate}[label=\arabic*.]
\item Norint nubrėžti tiesinės funkcijos grafiką, pakanka rasti du taškus, pažymėti juos koordinačių plokštumoje ir per juos nubrėžti tiese.
\item Norint rasti funkcijos nuli, sprendžiama lygtis $kx+b=0$.
\item Is lenteles atpazinkite tiesini sarysi: jei $x$ padidejus tuo paciu dydziu $y$ kinta tuo paciu dydziu, taskai yra vienoje tieseje.
\item Realaus uzdavinio formuleje $b$ daznai reiškia pastovia dali, o $k$ -- kaina, greiti arba pokyti vienam vienetui.
\end{enumerate}
\paragraph{Dažnos klaidos.}
Krypties koeficientas nėra tiesiog $y$ reikšmė. Jį reikia skaičiuoti kaip $\frac{\Delta y}{\Delta x}$. Braizant grafiką, asiu mastelis turi būti vienodas arba aiskiai pazymetas; kitaip grafikas gali klaidinti.

\subsection{Iliustruojantys pavyzdžiai}
\textbf{1 pavyzdys.} Raskime funkcijos $y=2x-3$ reikšmė, kai $x=4$.
\[
y=2\cdot4-3=5.
\]
\textbf{Atsakymas: $5$.}

\textbf{2 pavyzdys.} Raskime funkcijos $y=-3x+6$ nuli.
\[
-3x+6=0,\qquad x=2.
\]
\textbf{Atsakymas: $x=2$.}

\textbf{3 pavyzdys.} Raskime tieses, einancios per $A(1;2)$ ir $B(5;10)$, krypties koeficienta.
\[
k=\frac{10-2}{5-1}=2.
\]
\textbf{Atsakymas: $k=2$.}

\textbf{4 pavyzdys.} Taksi kaina sudaro $2$ eurai už isedima ir $0{,}80$ euro už kilometra. Parasykime kainos formule ir raskime kaina už $12$ km.
\[
y=0{,}8x+2,\qquad y=0{,}8\cdot12+2=11{,}6.
\]
\textbf{Atsakymas: $11{,}60$ euro.}

\paragraph{Pasitikrinkite.}
\begin{enumerate}[label=\alph*)]
\item Nubrezimui parinkite du funkcijos $y=x+4$ taškus.
\item Raskite funkcijos $y=5-2x$ nuli.
\item Apskaiciuokite krypties koeficienta per taškus $(2;7)$ ir $(6;15)$.
\end{enumerate}

\subsection{Ryšys su kitomis temomis}
Tiesinė funkcija sieja koordinačių plokštumą iš \ref{sec:dydziu-priklausomybes-ir-koordinaciu-plokstuma} su lygciu sprendimu \ref{sec:tiesines-lygtys-nelygybes-ir-sistemos}. Jį ruosia bendresne funkciju temą \ref{sec:funkcijos-ir-ju-savybes}, tiesinius ir netiesinius sarysius \ref{sec:tiesiniai-ir-netiesiniai-sarysiai-7-8} ir kvadratinę funkcija \ref{sec:kvadratine-funkcija}.

\section{Pitagoro teorema}\label{sec:pitos-teorema}

\subsection{Apibrėžtys}
\begin{apibreztis}
Statusis trikampis yra trikampis, kurio vienas kampas lygus $90^\circ$.
\end{apibreztis}
\begin{apibreztis}
Staciojo trikampio kraštinės, sudarancios statuji kampa, vadinamos statiniais, o pries statuji kampa esanti krastine -- įžambinę.
\end{apibreztis}
\begin{apibreztis}
Atvirkštinė teorema yra teiginys, gaunamas sukeičiant pradinės teoremos sąlygą ir išvadą. Jį ne visada teisinga, todėl turi būti irodyta atskirai.
\end{apibreztis}

\subsection{Teoremos ir savybės}
\begin{teorema}
Pitagoro teorema. Jei staciojo trikampio statiniai yra $a$ ir $b$, o įžambinę yra $c$, tai
\[
a^2+b^2=c^2.
\]
\end{teorema}
\begin{teorema}
Pitagoro teoremos atvirkstine teorema. Jei trikampio krastinems galioja $a^2+b^2=c^2$, kai $c$ yra ilgiausia krastine, tai trikampis statusis.
\end{teorema}
\begin{savybe}
Koordinačiu plokštumoje atstumas tarp tasku $A(x_1;y_1)$ ir $B(x_2;y_2)$ yra
\[
AB=\sqrt{(x_2-x_1)^2+(y_2-y_1)^2}.
\]
\end{savybe}
\begin{pastaba}
Pitagoro teorema taikoma tik staciajam trikampiui. Jei trikampis nezinomas, pirmiausia reikia irodyti arba patikrinti, kad jis statusis.
\end{pastaba}

\subsection{Taisyklės ir algoritmai}
\begin{enumerate}[label=\arabic*.]
\item Randant įžambinę, sudėkite statinių kvadratus ir ištraukite šaknį: $c=\sqrt{a^2+b^2}$.
\item Randant statini, iš izambines kvadrato atimkite kito statinio kvadrata: $a=\sqrt{c^2-b^2}$.
\item Tikrinant, ar trikampis statusis, ilgiausia krastine laikykite galima įžambinę ir tikrinkite kvadratu lygybe.
\item Koordinačiu plokštumoje nubreztas horizontalus ir vertikalus pokytis sudaro statinius.
\end{enumerate}
\paragraph{Dažnos klaidos.}
Izambine visada yra ilgiausia krastine ir yra pries statuji kampa. Negalima skaičiuoti $c=a+b$; Pitagoro teoremoje sudedami kraštinių kvadratai, o ne kraštinės.

\subsection{Iliustruojantys pavyzdžiai}
\textbf{1 pavyzdys.} Staciojo trikampio statiniai $6$ cm ir $8$ cm. Raskime įžambinę.
\[
c^2=6^2+8^2=36+64=100,\qquad c=10.
\]
\textbf{Atsakymas: $10$ cm.}

\textbf{2 pavyzdys.} Izambine $13$ cm, vienas statinis $5$ cm. Raskime kita statini.
\[
b^2=13^2-5^2=169-25=144,\qquad b=12.
\]
\textbf{Atsakymas: $12$ cm.}

\textbf{3 pavyzdys.} Ar trikampis su krastinemis $7$, $24$, $25$ yra statusis?
\[
7^2+24^2=49+576=625,\qquad 25^2=625.
\]
Lygybe teisinga, todėl trikampis statusis.

\textbf{4 pavyzdys.} Raskime atstuma tarp $A(-1;2)$ ir $B(5;10)$.
\[
AB=\sqrt{(5-(-1))^2+(10-2)^2}=\sqrt{6^2+8^2}=10.
\]
\textbf{Atsakymas: $10$.}

\paragraph{Pasitikrinkite.}
\begin{enumerate}[label=\alph*)]
\item Staciojo trikampio statiniai $9$ ir $12$. Raskite įžambinę.
\item Patikrinkite, ar kraštinės $8$, $15$, $17$ sudaro statuji trikampi.
\item Raskite atstuma tarp $(2;-3)$ ir $(6;0)$.
\end{enumerate}

\subsection{Ryšys su kitomis temomis}
Pitagoro teorema remiasi saknimis \ref{sec:saknys-ir-realiojo-skaiciaus-samprata} ir trikampiu savybemis \ref{sec:kampai-trikampiai-ir-keturkampiai}. Jį naudojama koordinačių geometrijoje, erdviniuose kunuose \ref{sec:erdviniai-kunai-ir-ju-pjuviai}, vektoriuose \ref{sec:mastelis-breziniai-ir-transformacijos} ir stereometrijoje \ref{sec:erdves-geometrija}.

\section{Panašūs trikampiai ir Talio teorema}\label{sec:panasus-trikampiai-ir-talio-teorema}

\subsection{Apibrėžtys}
\begin{apibreztis}
Du trikampiai yra panašus, jei ju atitinkami kampai lygus, o atitinkamu kraštinių ilgiai proporcingi.
\end{apibreztis}
\begin{apibreztis}
Panasumo koeficientas yra atitinkamu panašios figuros kraštinių ilgiu santykis. Jei iš mazesnes figuros pereiname i didesne, koeficientas daznai didesnis už $1$.
\end{apibreztis}
\begin{apibreztis}
Atitinkamos kraštinės yra kraštinės, esancios pries lygius kampus panasiose figurose.
\end{apibreztis}

\subsection{Teoremos ir savybės}
\begin{teorema}
Jei dvieju trikampių du atitinkami kampai lygus, tai trikampiai yra panašus.
\end{teorema}
\begin{teorema}
Talio teorema. Jei kelios lygiagrecios tieses kerta dvi tieses, tai jos atkerta proporcingas atkarpas.
\end{teorema}
\begin{savybe}
Panasiu figuru perimetrai santykiauja kaip atitinkamos kraštinės, o plotai -- kaip panasumo koeficiento kvadratas.
\end{savybe}

\subsection{Taisyklės ir algoritmai}
\begin{enumerate}[label=\arabic*.]
\item Pirmiausia nustatykite, kurie kampai atitinka vienas kita.
\item Sudarykite atitinkamu kraštinių proporcija ta pacia kryptimi: mazojo su mazojo, didziojo su didziojo.
\item Jei ieskote ilgio realiame kontekste, patikrinkite vienetus ir ar atsakymas logiskas.
\item Talio teoremoje pazymekite lygiagrecias tieses ir dvi kertamasias, tada užrašykite atitinkamu atkarpu santykius.
\end{enumerate}
\paragraph{Dažnos klaidos.}
Netinka proporcijoje maisyti neatitinkamų kraštinių. Taip pat reikia atskirti panašumą nuo lygumo: panašus trikampiai gali būti skirtingo dydžio, o lygus trikampiai yra ir panašus, ir vienodo dydžio.

\subsection{Iliustruojantys pavyzdžiai}
\textbf{1 pavyzdys.} Trikampiai panašus. Atitinkamos kraštinės yra $4$ cm ir $10$ cm. Mazojo trikampio kita krastine $6$ cm. Raskime atitinkama didziojo trikampio krastine.
\[
k=\frac{10}{4}=2{,}5,\qquad 6\cdot2{,}5=15.
\]
\textbf{Atsakymas: $15$ cm.}

\textbf{2 pavyzdys.} Trikampiai $ABC$ ir $DEF$ panašus, $AB=6$, $DE=9$, $AC=8$. Raskime $DF$, jei $AB$ atitinka $DE$, o $AC$ atitinka $DF$.
\[
\frac{AB}{DE}=\frac{AC}{DF},\qquad \frac69=\frac8{DF}.
\]
\[
6DF=72,\qquad DF=12.
\]
\textbf{Atsakymas: $12$.}

\textbf{3 pavyzdys.} Dvi lygiagrecios tieses atkerta ant vienos kertamosios atkarpas $3$ ir $5$, o ant kitos pirma atitinkama atkarpa yra $6$. Raskime antraja.
\[
\frac35=\frac6x,\qquad 3x=30,\qquad x=10.
\]
\textbf{Atsakymas: $10$.}

\paragraph{Pasitikrinkite.}
\begin{enumerate}[label=\alph*)]
\item Panasumo koeficientas iš mazos figuros i didele yra $3$. Mazoji krastine $7$ cm. Raskite atitinkama didziaja krastine.
\item Panasiu trikampių atitinkamos kraštinės $5$ ir $8$. Mazojo trikampio perimetras $30$. Raskite didziojo perimetra.
\item Jei panasumo koeficientas $4$, kiek kartu skiriasi plotai?
\end{enumerate}

\subsection{Ryšys su kitomis temomis}
Panasumas remiasi proporcijomis \ref{sec:procentai-proporcijos-ir-mastelis} ir kampu savybemis \ref{sec:kampai-trikampiai-ir-keturkampiai}. Jis naudojamas mastelyje \ref{sec:mastelis-breziniai-ir-transformacijos}, Pitagoro teoremos taikymuose \ref{sec:pitos-teorema} ir plotu bei turiu palyginimuose \ref{sec:plotai-turiai-ir-panasumas}.

\section{Transformacijos ir simetrijos}\label{sec:transformacijos-ir-simetrijos}

\subsection{Apibrėžtys}
\begin{apibreztis}
Plokstumos transformacija yra taisykle, kuri kiekvienam plokstumos taskui priskiria kita plokstumos taska.
\end{apibreztis}
\begin{apibreztis}
Asine simetrija yra transformacija tieses atzvilgiu: taskas ir jo vaizdas yra vienodų atstumu nuo simetrijos asies, o ju jungianti atkarpa statmena asiai.
\end{apibreztis}
\begin{apibreztis}
Centrine simetrija centro $O$ atzvilgiu yra transformacija, kai $O$ yra atkarpos $AA'$ vidurio taskas.
\end{apibreztis}
\begin{apibreztis}
Poslinkis perkelia visus figuros taškus ta pacia kryptimi ir tuo paciu atstumu.
\end{apibreztis}

\subsection{Teoremos ir savybės}
\begin{savybe}
Asine simetrija, centrine simetrija, poslinkis ir pasukimas issaugo atstumus, kampų didumus, plotus ir figuros forma.
\end{savybe}
\begin{savybe}
Koordinačiu plokštumoje atspindint taska $(x;y)$ $x$ asies atzvilgiu gaunama $(x;-y)$, $y$ asies atzvilgiu -- $(-x;y)$, o koordinačių pradzios atzvilgiu -- $(-x;-y)$.
\end{savybe}
\begin{savybe}
Jei poslinkio vektorius yra $(a;b)$, tai taskas $(x;y)$ pereina i $(x+a;y+b)$.
\end{savybe}

\subsection{Taisyklės ir algoritmai}
\begin{enumerate}[label=\arabic*.]
\item Atspindint taska tieses atzvilgiu, nubreziamas statmuo i asi ir toks pat atstumas atidedamas kitoje puseje.
\item Atliekant poslinki, kiekvienai virsunei pridedamas tas pats koordinačių pokytis.
\item Pagrindziant figuru lyguma, galima nurodyti transformaciju seka, kuri viena figura perkelia i kita.
\item Koordinačiu uzdaviniuose patogu sudaryti lentele: pradinis taskas, transformacijos taisykle, gautas taskas.
\end{enumerate}
\paragraph{Dažnos klaidos.}
Atspindint $x$ asies atzvilgiu keiciasi tik $y$ koordinates ženklas, o ne abi koordinates. Poslinkis nėra didinimas: figuros dydis nesikeicia.

\subsection{Iliustruojantys pavyzdžiai}
\textbf{1 pavyzdys.} Taskas $A(3;-4)$ atspindimas $x$ asies atzvilgiu.
\[
A'(3;4).
\]
\textbf{Atsakymas: $A'(3;4)$.}

\textbf{2 pavyzdys.} Taskui $B(-2;5)$ pritaikomas poslinkis $4$ vienetais i dešinę ir $3$ vienetais zemyn.
\[
x'=-2+4=2,\qquad y'=5-3=2.
\]
\textbf{Atsakymas: $B'(2;2)$.}

\textbf{3 pavyzdys.} Trikampio virsunes $A(0;1)$, $B(2;1)$, $C(1;4)$ pastumiamos vektoriumi $(3;-2)$. Raskime vaizdo virsunes.
\[
A'(3;-1),\qquad B'(5;-1),\qquad C'(4;2).
\]
\textbf{Atsakymas: $A'(3;-1)$, $B'(5;-1)$, $C'(4;2)$.}

\paragraph{Pasitikrinkite.}
\begin{enumerate}[label=\alph*)]
\item Atspindekite $P(-5;2)$ $y$ asies atzvilgiu.
\item Taska $Q(1;-3)$ pastumkite vektoriumi $(-4;6)$.
\item Paaiskinkite, kodel atspindeta figura yra lygi pradinei.
\end{enumerate}

\subsection{Ryšys su kitomis temomis}
Transformacijos pratęsia simetrijos pradžią \ref{sec:simetrija-ir-transformacijos-pradzia} ir koordinačių plokštuma \ref{sec:dydziu-priklausomybes-ir-koordinaciu-plokstuma}. Jos padeda suprasti figuru lyguma, panašumą \ref{sec:panasus-trikampiai-ir-talio-teorema}, vektorius ir breziniu skaityma \ref{sec:mastelis-breziniai-ir-transformacijos}.

\section{Erdviniai kūnai ir jų pjūviai}\label{sec:erdviniai-kunai-ir-ju-pjuviai}

\subsection{Apibrėžtys}
\begin{apibreztis}
Erdvinis kunas yra trimatis objektas, turintis ilgį, ploti ir auksti.
\end{apibreztis}
\begin{apibreztis}
Prizme yra briaunainis, kurio du pagrindai yra lygus ir lygiagretus daugiakampiai, o sonines sienos -- lygiagretainiai. Staciojoje prizmeje sonines briaunos statmenos pagrindams.
\end{apibreztis}
\begin{apibreztis}
Piramide yra briaunainis, kurio vienas pagrindas yra daugiakampis, o sonines sienos -- trikampiai, sueinantys i viena virsune.
\end{apibreztis}
\begin{apibreztis}
Pjuvis yra plokstumos ir erdvinio kuno bendroji dalis.
\end{apibreztis}

\subsection{Teoremos ir savybės}
\begin{savybe}
Staciosios prizmes tūris lygus pagrindo ploto ir aukscio sandaugai:
\[
V=S_p h.
\]
\end{savybe}
\begin{savybe}
Piramides tūris apskaiciuojamas formule
\[
V=\frac13 S_p h.
\]
\end{savybe}
\begin{savybe}
Ritinio tūris ir soninio pavirsiaus plotas:
\[
V=\pi r^2h,\qquad S_{\text{son}}=2\pi rh.
\]
Kugio tūris:
\[
V=\frac13\pi r^2h.
\]
\end{savybe}
\begin{savybe}
Staciakampio gretasienio erdvine istrizaine tenkina
\[
d^2=a^2+b^2+c^2.
\]
\end{savybe}

\subsection{Taisyklės ir algoritmai}
\begin{enumerate}[label=\arabic*.]
\item Pirmiausia atpazinkite kuna ir pazymekite, kas yra pagrindas, aukstine, spindulys, sudaromoji ar apotema.
\item Pavirsiaus plotui rasti sudėkite visu sienu arba daliu plotus.
\item Turiui rasti naudokite tinkama formule ir butinai vienodus matavimo vienetus.
\item Tiriant pjuvi, pazymekite plokstumos susikirtimo taškus su briaunomis ir junkite taškus tose paciose sienose.
\end{enumerate}
\paragraph{Dažnos klaidos.}
Pavirsiaus plotas matuojamas kvadratiniais vienetais, o tūris -- kubiniais. Ritinio ir kugio formulese spindulys nėra skersmuo; jei duotas skersmuo, ji reikia padalyti iš $2$.

\subsection{Iliustruojantys pavyzdžiai}
\textbf{1 pavyzdys.} Staciosios prizmes pagrindo plotas $18\ \text{cm}^2$, aukstis $7$ cm. Raskime turi.
\[
V=S_p h=18\cdot7=126.
\]
\textbf{Atsakymas: $126\ \text{cm}^3$.}

\textbf{2 pavyzdys.} Staciakampio gretasienio matmenys $3$, $4$ ir $12$ cm. Raskime erdvine istrizaine.
\[
d^2=3^2+4^2+12^2=169,\qquad d=13.
\]
\textbf{Atsakymas: $13$ cm.}

\textbf{3 pavyzdys.} Ritinio spindulys $4$ cm, aukstis $10$ cm. Raskime turi.
\[
V=\pi r^2h=\pi\cdot4^2\cdot10=160\pi.
\]
\textbf{Atsakymas: $160\pi\ \text{cm}^3$.}

\textbf{4 pavyzdys.} Taisyklingosios piramides pagrindo plotas $48\ \text{cm}^2$, aukstis $9$ cm. Raskime turi.
\[
V=\frac13\cdot48\cdot9=144.
\]
\textbf{Atsakymas: $144\ \text{cm}^3$.}

\paragraph{Pasitikrinkite.}
\begin{enumerate}[label=\alph*)]
\item Kubo briauna $6$ cm. Raskite turi.
\item Ritinio skersmuo $10$ cm, aukstis $8$ cm. Raskite turi.
\item Paaiskinkite, kodel tas pats kubas gali tureti skirtingos formos pjuvius.
\end{enumerate}

\subsection{Ryšys su kitomis temomis}
Erdviniai kunai remiasi ploto, pavirsiaus ir turio ziniomis \ref{sec:plotas-pavirsius-ir-turis}. Ju uzdaviniuose taikoma Pitagoro teorema \ref{sec:pitos-teorema}, panašus trikampiai \ref{sec:panasus-trikampiai-ir-talio-teorema}, apskritimo ir skritulio formulės bei velesne erdves geometrija \ref{sec:erdves-geometrija}.

\section{Statistika, imtys ir sklaida}\label{sec:statistika-imtys-ir-sklaida}

\subsection{Apibrėžtys}
\begin{apibreztis}
Populiacija yra visa tiriamu objektu grupe, o imtis -- iš populiacijos atrinkta dalis.
\end{apibreztis}
\begin{apibreztis}
Imties dydis yra imties elementu skaičius. Reprezentatyvioji imtis pakankamai gerai atspindi populiacija.
\end{apibreztis}
\begin{apibreztis}
Aritmetinis vidurkis yra visu duomenų reikšmių suma, padalyta iš reikšmių skaičiaus. Mediana yra surikiuotos duomenų eiles vidurine reikšmė, o moda -- dazniausiai pasikartojanti reikšmė.
\end{apibreztis}
\begin{apibreztis}
Sklaidos plotis yra didziausios ir maziausios reikšmių skirtumas. Kvartiliai $Q_1$, $Q_2$, $Q_3$ dalija surikiuotus duomenis i keturias mazdaug lygias dalis.
\end{apibreztis}

\subsection{Teoremos ir savybės}
\begin{savybe}
Vidurkis jautrus isskirtims, o mediana ir kvartiliai paprastai geriau apibudina duomenų centra, kai yra labai dideliu arba labai mažų reikšmių.
\end{savybe}
\begin{savybe}
Dazniu lenteleje bendra duomenų suma gaunama sudedant kiekvienos reikšmes ir jos daznio sandaugas.
\end{savybe}
\begin{savybe}
Santykinis daznis lygus $\frac{\text{daznis}}{\text{visu stebejimu skaičius}}$, o sukauptasis daznis rodo, kiek reikšmių nevirsija pasirinktos ribos.
\end{savybe}
\begin{pastaba}
Statistine išvadą patikimesne, kai imtis sudaryta atsitiktinai, yra pakankamai didele ir nepalanki jokiai iš anksto pasirinktai isvadai.
\end{pastaba}

\subsection{Taisyklės ir algoritmai}
\begin{enumerate}[label=\arabic*.]
\item Vidurkiui rasti sudėkite visas reikšmes ir padalykite iš ju skaičiaus.
\item Medianai rasti duomenis surikiuokite. Jei reikšmių skaičius nelyginis, imkite vidurine; jei lyginis, imkite dvieju viduriniu vidurki.
\item Sklaidos plociui rasti iš didziausios reikšmes atimkite mažiausia.
\item Skritulineje diagramoje sektoriaus kampas lygus santykiniam dazniui, padaugintam iš $360^\circ$.
\item Staciakampes diagramos su usais santrauka: minimumas, $Q_1$, mediana, $Q_3$, maksimumas.
\end{enumerate}
\paragraph{Dažnos klaidos.}
Pries ieskant medianos duomenis butina surikiuoti. Skritulineje diagramoje procentai turi sudaryti $100\%$. Imtis nėra tas pats, kas populiacija: iš mazos ar saliskos imties negalima drasiai spresti apie visus.

\subsection{Iliustruojantys pavyzdžiai}
\textbf{1 pavyzdys.} Raskime duomenų $4,7,7,10,12$ vidurki.
\[
\overline{x}=\frac{4+7+7+10+12}{5}=8.
\]
\textbf{Atsakymas: $8$.}

\textbf{2 pavyzdys.} Raskime duomenų $9,3,12,5,5,8$ mediana.
Surikiuojame:
\[
3,5,5,8,9,12.
\]
Vidurines reikšmes $5$ ir $8$, todėl
\[
\text{mediana}=\frac{5+8}{2}=6{,}5.
\]
\textbf{Atsakymas: $6{,}5$.}

\textbf{3 pavyzdys.} Klaseje apklausti $24$ mokiniai: $6$ pasirinko krepsini, $10$ -- futbola, $8$ -- plaukima. Koks futbolo sektoriaus kampas skritulineje diagramoje?
\[
\frac{10}{24}\cdot360^\circ=150^\circ.
\]
\textbf{Atsakymas: $150^\circ$.}

\textbf{4 pavyzdys.} Duomenys: $2,4,5,7,9,10,12,15$. Raskime kvartilius.
Mediana:
\[
Q_2=\frac{7+9}{2}=8.
\]
Apatines dalies $2,4,5,7$ mediana:
\[
Q_1=\frac{4+5}{2}=4{,}5.
\]
Viršutines dalies $9,10,12,15$ mediana:
\[
Q_3=\frac{10+12}{2}=11.
\]
\textbf{Atsakymas: $Q_1=4{,}5$, $Q_2=8$, $Q_3=11$.}

\paragraph{Pasitikrinkite.}
\begin{enumerate}[label=\alph*)]
\item Raskite duomenų $6,6,9,11,13$ vidurki ir mediana.
\item Duomenų rinkinyje $18,20,21,22,60$ paaiškinkite, kuri reikšmė yra išskirtis.
\item Jei kategorija sudaro $35\%$, koks jos sektoriaus kampas skritulineje diagramoje?
\end{enumerate}

\subsection{Ryšys su kitomis temomis}
Statistika tęsia duomenų lentelių ir diagramų temas \ref{sec:duomenu-lenteles-ir-diagramos} bei vidurkius \ref{sec:statistiniai-duomenys-ir-vidurkiai}. Jį siejasi su procentais \ref{sec:procentai-proporcijos-ir-mastelis}, tikimybė \ref{sec:kombinatorikos-pradzia-ir-tikimybe} ir velesne koreliacija \ref{sec:statistiniai-rodikliai-ir-koreliacija}.

\section{Kombinatorikos pradžia ir tikimybė}\label{sec:kombinatorikos-pradzia-ir-tikimybe}

\subsection{Apibrėžtys}
\begin{apibreztis}
Baigtinis bandymas yra veiksmas arba situacija, turinti baigtini galimu baigčių skaičių.
\end{apibreztis}
\begin{apibreztis}
Ivykis yra bandymo baigčių aibė. Palankios baigtys yra tos, kurioms įvykis ivyksta.
\end{apibreztis}
\begin{apibreztis}
Klasikine tikimybė, kai visos baigtys vienodai galimos, skaičiuojama
\[
P(A)=\frac{\text{palankiu baigčių skaičius}}{\text{visu baigčių skaičius}}.
\]
\end{apibreztis}
\begin{apibreztis}
Daugybos taisykle: jei pirma pasirinkima galima atlikti $m$ budais, o po jo antra -- $n$ budais, tai abu pasirinkimus galima atlikti $mn$ budu.
\end{apibreztis}

\subsection{Teoremos ir savybės}
\begin{savybe}
Bet kurio įvykio tikimybė tenkina
\[
0\le P(A)\le1.
\]
\end{savybe}
\begin{savybe}
Neimanomo įvykio tikimybė lygi $0$, o butino įvykio tikimybė lygi $1$.
\end{savybe}
\begin{savybe}
Jei ivykiai nesikerta, ju sajungos tikimybė lygi tikimybių sumai.
\end{savybe}
\begin{savybe}
Priesingo įvykio tikimybė:
\[
P(\overline{A})=1-P(A).
\]
\end{savybe}

\subsection{Taisyklės ir algoritmai}
\begin{enumerate}[label=\arabic*.]
\item Apskaiciuojant klasikine tikimybė, isvardykite visas vienodai galimas baigtis.
\item Suskaiciuokite palankias baigtis ir sudarykite trupmena.
\item Jei pasirinkimai atliekami vienas po kito, naudokite daugybos taisykle.
\item Jei patogiau skaičiuoti, raskite priesingo įvykio tikimybė ir atimkite iš $1$.
\end{enumerate}
\paragraph{Dažnos klaidos.}
Tikimybė negali būti didesne už $1$ ar neigiama. Ne visi bandymai yra klasikiniai: jei baigtys nevienodai galimos, vien palankiu baigčių skaiciavimas gali klaidinti.

\subsection{Iliustruojantys pavyzdžiai}
\textbf{1 pavyzdys.} Metamas taisyklingas kauliukas. Raskime tikimybė ismesti lygini skaičių.
Palankios baigtys $2,4,6$, ju yra $3$, iš viso baigčių $6$:
\[
P=\frac36=\frac12.
\]
\textbf{Atsakymas: $\frac12$.}

\textbf{2 pavyzdys.} Mokinys renkasi vienus marskinelius iš $4$ ir vienas kelnes iš $3$. Kiek deriniu galima sudaryti?
\[
4\cdot3=12.
\]
\textbf{Atsakymas: $12$ deriniu.}

\textbf{3 pavyzdys.} Dezeje yra $5$ raudoni ir $7$ melyni rutuliai. Atsitiktinai traukiamas vienas rutulys. Raskime tikimybė istraukti raudona.
\[
P=\frac{5}{5+7}=\frac5{12}.
\]
\textbf{Atsakymas: $\frac5{12}$.}

\textbf{4 pavyzdys.} Is skaičių $1,2,3,\ldots,20$ atsitiktinai pasirenkamas vienas. Kokia tikimybė, kad jis nesidalija iš $5$?
Is $5$ dalijasi $5,10,15,20$, tai $4$ skaičiai. Priesingo įvykio tikimybė $\frac4{20}=\frac15$, todėl
\[
P=1-\frac15=\frac45.
\]
\textbf{Atsakymas: $\frac45$.}

\paragraph{Pasitikrinkite.}
\begin{enumerate}[label=\alph*)]
\item Metama moneta ir kauliukas. Kiek yra visu baigčių?
\item Kokia tikimybė kauliuku ismesti skaičių, didesni už $4$?
\item Slaptazodi sudaro viena raide iš $5$ galimu ir vienas skaitmuo iš $10$ galimu. Kiek slaptazodziu galima sudaryti?
\end{enumerate}

\subsection{Ryšys su kitomis temomis}
Kombinatorika ir tikimybė pratęsia atsitiktinių bandymų temą \ref{sec:atsitiktiniai-bandymas-ir-tikimybe}. Jos remiasi trupmenomis \ref{sec:paprastosios-ir-desimtaines-trupmenos}, procentais \ref{sec:procentai-proporcijos-ir-mastelis}, statistika \ref{sec:statistika-imtys-ir-sklaida} ir ruosia tikimybes bei kombinatorikos temą \ref{sec:tikimybe-ir-kombinatorika}.

\section{Finansiniai skaičiavimai: palūkanos, kreditas ir lizingas}\label{sec:finansiniai-skaiciavimai-palukanos-kreditas-lizingas}

\subsection{Apibrėžtys}
\begin{apibreztis}
Palukanos yra pinigu suma, mokama už pasiskolintus pinigus arba gaunama už laikomus pinigus. Palukanu norma nurodo, kokia pradinės sumos dalis priskaiciuojama per tam tikra laika.
\end{apibreztis}
\begin{apibreztis}
Paprastosios palukanos skaiciuojamos tik nuo pradinės sumos:
\[
S=P(1+rt),
\]
kur $P$ -- pradine suma, $r$ -- laikotarpio palūkanų norma desimtaine išraiška, $t$ -- laikotarpiu skaičius.
\end{apibreztis}
\begin{apibreztis}
Sudėtines palukanos skaiciuojamos nuo vis augancios sumos:
\[
S=P(1+r)^t.
\]
\end{apibreztis}
\begin{apibreztis}
Kreditas yra skolinimasis, kai pinigai grazinami veliau su sutartomis palukanomis ir mokesciais. Lizingas yra daikto isigijimas mokant dalimis, kol ivykdomos sutarties salygos.
\end{apibreztis}

\subsection{Teoremos ir savybės}
\begin{savybe}
Bendra finansinio pasiulymo kaina priklauso ne tik nuo palūkanų normos, bet ir nuo sutarties mokesciu, termino, imoku skaičiaus, pradinio inaso ir palūkanų skaičiavimo budo.
\end{savybe}
\begin{savybe}
Ta pati procentine norma ilgesniam laikui lemia didesne palūkanų suma. Sudetiniu palūkanų atveju augimas greiteja, nes palukanos skaiciuojamos ir nuo anksciau priskaiciuotu palūkanų.
\end{savybe}
\begin{pastaba}
Finansiniuose sprendimuose matematika padeda palyginti variantus, bet reikia vertinti ir rizika: ar imoka bus pakeliama, ar palūkanų norma gali keistis, ar yra papildomu salygu.
\end{pastaba}

\subsection{Taisyklės ir algoritmai}
\begin{enumerate}[label=\arabic*.]
\item Procentu norma paverskite desimtainiu skaičiumi: $12\%=0{,}12$.
\item Atskirkite pradine suma, laikotarpiu skaičių, palūkanų norma ir papildomus mokescius.
\item Paprastosioms palukanoms taikykite $S=P(1+rt)$, sudetinems -- $S=P(1+r)^t$.
\item Lygindami kreditus ar lizinga, skaiciuokite bendra sumokama suma: pradinis inasas + visos imokos + mokesciai.
\item Biudzete pajamas ir islaidas rasykite atskirai, o likuti skaiciuokite kaip pajamas minus islaidas.
\end{enumerate}
\paragraph{Dažnos klaidos.}
Metinė palūkanų norma negali būti tiesiogiai taikoma mėnesiui, jei uzdavinys prašo mėnesinio skaičiavimo: reikia aiskiai zinoti laikotarpi. Taip pat mažiausia menesine imoka ne visada reiškia mažiausia bendra kaina.

\subsection{Iliustruojantys pavyzdžiai}
\textbf{1 pavyzdys.} I banko saskaita padeta $500$ euru, metine paprastuju palūkanų norma $4\%$, laikotarpis $3$ metai. Raskime sukaupta suma.
\[
S=500(1+0{,}04\cdot3)=500\cdot1{,}12=560.
\]
\textbf{Atsakymas: $560$ euru.}

\textbf{2 pavyzdys.} Ta pati $500$ euru suma kaupiama su $4\%$ metinemis sudetinemis palukanomis $3$ metus.
\[
S=500(1{,}04)^3=500\cdot1{,}124864=562{,}432.
\]
\textbf{Atsakymas: apie $562{,}43$ euro.}

\textbf{3 pavyzdys.} Telefonas kainuoja $360$ euru. Perkant lizingu mokamas $40$ euru pradinis inasas ir $12$ menesiniu imoku po $29$ eurus. Kiek permokama?
\[
\text{bendra suma}=40+12\cdot29=388.
\]
\[
388-360=28.
\]
\textbf{Atsakymas: permokama $28$ eurai.}

\textbf{4 pavyzdys.} Mokinys planuoja rengini. Pajamos $320$ euru, sales nuoma $120$ euru, maistas $95$ eurai, reklama $36$ eurai. Koks likutis?
\[
320-(120+95+36)=320-251=69.
\]
\textbf{Atsakymas: $69$ eurai.}

\paragraph{Pasitikrinkite.}
\begin{enumerate}[label=\alph*)]
\item Apskaiciuokite $800$ euru suma po $2$ metu su $5\%$ paprastosiomis palukanomis.
\item Preke kainuoja $240$ euru, lizingo bendra suma $267$ eurai. Kiek procentu permokama nuo pradinės kainos?
\item Biudzeto pajamos $150$ euru, islaidos $47$, $38$ ir $29$ eurai. Raskite likuti.
\end{enumerate}

\subsection{Ryšys su kitomis temomis}
Finansiniai skaičiavimai remiasi procentais \ref{sec:procentai-proporcijos-ir-mastelis}, laipsniais \ref{sec:racionalieji-skaiciai-ir-laipsniai}, tiesiniais modeliais \ref{sec:tiesine-funkcija-ir-grafikas} ir duomenų palyginimu \ref{sec:diagramu-analize-ir-sklaida}. Jie moko taikyti matematika realiems pasirinkimams.

\section{Algebrinių reiškinių pertvarkiai ir skaidymas}\label{sec:algebriniu-reiskiniu-pertvarkiai-ir-skaidymas}

\subsection{Apibrėžtys}
\begin{apibreztis}
Tapatybiskai lygus reiškiniai yra reiškiniai, kuriu reikšmes sutampa visoms leistinoms kintamųjų reiksmems.
\end{apibreztis}
\begin{apibreztis}
Skaidyti dauginamaisiais reiškia užrašyti reiškinį kaip keliu reiškinių sandauga.
\end{apibreztis}
\begin{apibreztis}
Bendras daugiklis yra daugiklis, kuris yra kiekviename daugianario naryje.
\end{apibreztis}
\begin{apibreztis}
Dvinario kvadrato isskyrimas -- kvadratinio trinaro pertvarkymas i pavidala $(x+a)^2+b$ arba $k(x+a)^2+b$.
\end{apibreztis}

\subsection{Teoremos ir savybės}
\begin{savybe}
Jei kiekvienas daugianario narys turi bendra daugikli $d$, tai
\[
da+db=d(a+b).
\]
\end{savybe}
\begin{savybe}
Dažnosios daugybos formulės gali būti taikomos abiem kryptimis:
\[
a^2-b^2=(a-b)(a+b),
\]
\[
a^2+2ab+b^2=(a+b)^2,\qquad a^2-2ab+b^2=(a-b)^2.
\]
\end{savybe}
\begin{savybe}
Grupavimas leidzia skaidyti, kai keli nariai turi viena bendra daugikli, o kiti -- kita, bet po iskelimo lieka vienodas skliaustas.
\end{savybe}

\subsection{Taisyklės ir algoritmai}
\begin{enumerate}[label=\arabic*.]
\item Pirmiausia ieskokite didziausio bendro skaitinio ir raidinio daugiklio.
\item Patikrinkite, ar reiškinys yra kvadratu skirtumas arba pilnas dvinario kvadratas.
\item Jei nariu keturi, bandykite grupuoti po du.
\item Po skaidymo patikrinkite atsakymą isskleisdami sandauga atgal.
\end{enumerate}
\paragraph{Dažnos klaidos.}
Iskeliant neigiama bendra daugikli, skliaustuose keiciasi visi zenklai. Skaidant $a^2+b^2$ negalima naudoti kvadratu skirtumo formulės, nes tai nėra skirtumas.

\subsection{Iliustruojantys pavyzdžiai}
\textbf{1 pavyzdys.} Isskaidykime
\[
6x^2-9x.
\]
Bendras daugiklis $3x$:
\[
6x^2-9x=3x(2x-3).
\]
\textbf{Atsakymas: $3x(2x-3)$.}

\textbf{2 pavyzdys.} Isskaidykime
\[
x^2-25.
\]
Tai kvadratu skirtumas:
\[
x^2-25=x^2-5^2=(x-5)(x+5).
\]
\textbf{Atsakymas: $(x-5)(x+5)$.}

\textbf{3 pavyzdys.} Isskaidykime grupavimu:
\[
ax+ay+3x+3y.
\]
\[
ax+ay+3x+3y=a(x+y)+3(x+y)=(x+y)(a+3).
\]
\textbf{Atsakymas: $(x+y)(a+3)$.}

\textbf{4 pavyzdys.} Isskirkime dvinario kvadrata:
\[
x^2+6x+11.
\]
\[
x^2+6x+11=x^2+6x+9+2=(x+3)^2+2.
\]
\textbf{Atsakymas: $(x+3)^2+2$.}

\paragraph{Pasitikrinkite.}
\begin{enumerate}[label=\alph*)]
\item Isskaidykite $12a^2b-18ab^2$.
\item Isskaidykite $9x^2-16$.
\item Grupavimu išskaidykite $mx+my+2x+2y$.
\end{enumerate}

\subsection{Ryšys su kitomis temomis}
Skaidymas remiasi algebriniais reiskiniais \ref{sec:algebriniai-reiskiniai-ir-daugianariai}. Jis naudojamas lygciu sprendimui, trupmeniniu reiskiniu prastinimui, kvadratinems lygtims \ref{sec:kvadratines-lygtys-ir-nelygybes} ir funkciju tyrimui.

\section{Tiesinės nelygybės ir sprendinių aibės}\label{sec:tiesines-nelygybes-ir-sprendiniu-aibes}

\subsection{Apibrėžtys}
\begin{apibreztis}
Nelygybe su vienu nežinomuoju yra teiginys su zenklu $<$, $>$, $\le$ arba $\ge$, kuriame yra kintamasis.
\end{apibreztis}
\begin{apibreztis}
Nelygybes sprendinys yra kintamojo reikšmė, kuriai nelygybė teisinga. Visu sprendiniu rinkinys vadinamas sprendiniu aibė.
\end{apibreztis}
\begin{apibreztis}
Intervalas yra realiuju skaičių aibė, aprasyta nelygybė. Pavyzdziui, $(-2;5]=\{x\in\mathbb{R}:-2<x\le5\}$.
\end{apibreztis}
\begin{apibreztis}
Dviguboji nelygybė, pavyzdžiui, $2<x\le7$, reiškia dvi salygas kartu: $x>2$ ir $x\le7$.
\end{apibreztis}

\subsection{Teoremos ir savybės}
\begin{savybe}
Jei prie abieju nelygybės pusiu pridedamas arba atimamas tas pats skaičius, nelygybės ženklas nesikeicia.
\end{savybe}
\begin{savybe}
Dauginant arba dalijant nelygybė iš teigiamo skaičiaus, ženklas nesikeicia; dauginant arba dalijant iš neigiamo skaičiaus, ženklas apsivercia.
\end{savybe}
\begin{savybe}
Nelygybiu sistemos sprendiniu aibė yra visu jos nelygybiu sprendiniu aibiu sankirta.
\end{savybe}

\subsection{Taisyklės ir algoritmai}
\begin{enumerate}[label=\arabic*.]
\item Sprendziant nelygybė, elkitės kaip su lygtimi, bet pazymekite kiekviena dalyba ar daugyba iš neigiamo skaičiaus.
\item Sprendini skaičiu tiesėje: tuscias taskas reiškia $<$ arba $>$, uzpildytas taskas -- $\le$ arba $\ge$.
\item Dvigubaja nelygybė galima spręsti vienu metu visose trijose dalyse arba isskaidyti i sistema.
\item Sistemoje raskite bendra visu salygu dali.
\end{enumerate}
\paragraph{Dažnos klaidos.}
Zenklai $<$ ir $\le$ skiriasi: vienu atveju ribinis skaičius nepriklauso sprendiniu aibei, kitu priklauso. Intervalo uzrase skliaustas $($ arba $)$ reiškia, kad galas neitraukiamas, o $[$ arba $]$ -- itraukiamas.

\subsection{Iliustruojantys pavyzdžiai}
\textbf{1 pavyzdys.} Ispręskime
\[
4x-7\le9.
\]
\[
4x\le16,\qquad x\le4.
\]
\textbf{Atsakymas: $(-\infty;4]$.}

\textbf{2 pavyzdys.} Ispręskime
\[
-3x+2<11.
\]
\[
-3x<9.
\]
Dalyba iš $-3$ keicia ženklą:
\[
x>-3.
\]
\textbf{Atsakymas: $(-3;+\infty)$.}

\textbf{3 pavyzdys.} Ispręskime dvigubaja nelygybė
\[
-2\le 3x+1<10.
\]
Atimame $1$ iš visu daliu:
\[
-3\le3x<9.
\]
Daliname iš $3$:
\[
-1\le x<3.
\]
\textbf{Atsakymas: $[-1;3)$.}

\textbf{4 pavyzdys.} Raskime sistemos sprendinius:
\[
\begin{cases}
x>1,\\
x\le5.
\end{cases}
\]
Bendra dalis yra $1<x\le5$.
\textbf{Atsakymas: $(1;5]$.}

\paragraph{Pasitikrinkite.}
\begin{enumerate}[label=\alph*)]
\item Ispręskite $2x+5>17$.
\item Ispręskite $-5x\le20$.
\item Uzrasykite intervalu sprendinius $-4<x\le2$.
\end{enumerate}

\subsection{Ryšys su kitomis temomis}
Tiesines nelygybes remiasi lygciu sprendimu \ref{sec:tiesines-lygtys-nelygybes-ir-sistemos} ir realiuju skaiciu tvarka. Jos reikalingos intervalams \ref{sec:realieji-skaiciai-ir-intervalai}, funkciju reiksmiu tyrimui ir uzdaviniams su salygomis.

\section{Tiesiniai ir netiesiniai sąryšiai}\label{sec:tiesiniai-ir-netiesiniai-sarysiai-7-8}

\subsection{Apibrėžtys}
\begin{apibreztis}
Tiesioginis proporcingumas yra sarysis $y=kx$, kai santykis $\frac{y}{x}$ yra pastovus.
\end{apibreztis}
\begin{apibreztis}
Atvirkstinis proporcingumas yra sarysis
\[
y=\frac{k}{x},\qquad x\ne0,
\]
kai sandauga $xy$ yra pastovi.
\end{apibreztis}
\begin{apibreztis}
Tiesinis sarysis yra sarysis $y=kx+b$. Jis gali tureti pastovia dali $b$ ir pastovu pokyti $k$.
\end{apibreztis}
\begin{apibreztis}
Netiesinis sarysis yra priklausomybe, kurios grafikas nėra tiese, pavyzdžiui, atvirkstinio proporcingumo grafikas.
\end{apibreztis}

\subsection{Teoremos ir savybės}
\begin{savybe}
Tiesioginio proporcingumo grafikas yra tiese, einanti per koordinačių pradžią.
\end{savybe}
\begin{savybe}
Atvirkstinio proporcingumo atveju, kai vienas dydis padideja kelis kartus, kitas tiek pat kartu sumazeja, jei sandauga turi likti pastovi.
\end{savybe}
\begin{savybe}
Tiesinio saryšio lenteleje vienodi $x$ pokyčiai atitinka vienodus $y$ pokyčius; atvirkstinio proporcingumo lenteleje pastovi yra sandauga $xy$.
\end{savybe}

\subsection{Taisyklės ir algoritmai}
\begin{enumerate}[label=\arabic*.]
\item Is lenteles raskite, kas pastovu: santykis $\frac{y}{x}$, sandauga $xy$ ar pokytis $\frac{\Delta y}{\Delta x}$.
\item Jei pastovus santykis, rasykite $y=kx$.
\item Jei pastovi sandauga, rasykite $y=\frac{k}{x}$.
\item Jei pastovus pokytis, bet grafikas nebutinai eina per pradžią, rasykite $y=kx+b$.
\item Realaus konteksto atsakymą tikrinkite pagal prasmę: laikas, ilgis, kaina negali būti neigiami, jei situacija to neleidzia.
\end{enumerate}
\paragraph{Dažnos klaidos.}
Tiesioginis ir atvirkštinis proporcingumas abu gali būti aprašyti lentele, bet ju taisyklės skirtingos. Jei $x$ padvigubejus $y$ irgi padvigubeja, tai tiesioginis proporcingumas; jei $y$ sumazeja perpus, tai atvirkštinis.

\subsection{Iliustruojantys pavyzdžiai}
\textbf{1 pavyzdys.} Kelio ilgis pastovus $120$ km. Sudarykime greicio $v$ ir laiko $t$ sarysi.
\[
vt=120,\qquad t=\frac{120}{v}.
\]
Tai atvirkštinis proporcingumas.

\textbf{2 pavyzdys.} Darbo uzmokestis yra $6$ eurai už valanda. Jei dirbta $x$ valandu, uzmokestis $y=6x$. Tai tiesioginis proporcingumas.

\textbf{3 pavyzdys.} Sporto klubo abonementas kainuoja $15$ euru registracijos mokesti ir $8$ eurus už kiekviena apsilankyma. Formule:
\[
y=8x+15.
\]
Kai $x=5$, $y=8\cdot5+15=55$.
\textbf{Atsakymas: $55$ eurai.}

\textbf{4 pavyzdys.} Lenteleje pateikti taskai $(2;18)$, $(3;12)$, $(6;6)$. Nustatykime sarysi.
\[
2\cdot18=36,\quad 3\cdot12=36,\quad 6\cdot6=36.
\]
Sandauga pastovi, todėl
\[
y=\frac{36}{x}.
\]

\paragraph{Pasitikrinkite.}
\begin{enumerate}[label=\alph*)]
\item Jei $y=4x$, raskite $y$, kai $x=7$.
\item Jei $xy=60$, raskite $y$, kai $x=12$.
\item Lenteleje $x$: $1,2,3$; $y$: $5,8,11$. Koks sarysis?
\end{enumerate}

\subsection{Ryšys su kitomis temomis}
Sąryšiai remiasi proporcijomis \ref{sec:procentai-proporcijos-ir-mastelis}, koordinačių plokštuma ir tiesine funkcija \ref{sec:tiesine-funkcija-ir-grafikas}. Jie padeda modeliuoti greiti, kaina, darba, plota ir ruosia funkciju savybiu temą \ref{sec:funkcijos-ir-ju-savybes}.

\section{Mastelis, brėžiniai ir transformacijos}\label{sec:mastelis-breziniai-ir-transformacijos}

\subsection{Apibrėžtys}
\begin{apibreztis}
Brezinys yra tikslus geometrinis vaizdas, kuriame laikomasi sutartu zymejimu, matmenu ir mastelio.
\end{apibreztis}
\begin{apibreztis}
Mastelis $m:n$ rodo brezinio ir tikrojo dydžio santyki. Dazniausiai $1:n$ reiškia sumazinima, o $n:1$ -- padidinima.
\end{apibreztis}
\begin{apibreztis}
Vektorius yra kryptine atkarpa. Jis turi ilgį ir krypti. Lygus vektoriai turi vienoda ilgį ir krypti, nors gali prasideti skirtinguose taskuose.
\end{apibreztis}
\begin{apibreztis}
Objekto projekcija yra jo vaizdas iš tam tikros krypties: iš virsaus, iš priekio arba iš sono.
\end{apibreztis}

\subsection{Teoremos ir savybės}
\begin{savybe}
Jei planas nubraizytas masteliu $1:n$, tikrasis ilgis yra brezinio ilgis, padaugintas iš $n$.
\end{savybe}
\begin{savybe}
Vektoriu sudėtis atitinka nuoseklu poslinki: jei pirma pastumiama vektoriumi $\vec a$, o tada $\vec b$, bendras poslinkis yra $\vec a+\vec b$.
\end{savybe}
\begin{savybe}
Koordinačiu plokštumoje vektoriaus nuo $A(x_1;y_1)$ iki $B(x_2;y_2)$ koordinates yra
\[
\overrightarrow{AB}=(x_2-x_1;\ y_2-y_1).
\]
\end{savybe}

\subsection{Taisyklės ir algoritmai}
\begin{enumerate}[label=\arabic*.]
\item Mastelio uzdavinyje pirmiausia visus ilgius paverskite tais paciais vienetais.
\item Jei duotas brezinio ilgis ir mastelis $1:n$, tikraji ilgį dauginkite iš $n$; jei duotas tikrasis ilgis, brezinio ilgį dalykite iš $n$.
\item Vektoriaus koordinates raskite iš pabaigos koordinačių atimdami pradzios koordinates.
\item Projektuojant kuna, atskirai apgalvokite, kurie matmenys matomi iš virsaus, priekio ir sono.
\end{enumerate}
\paragraph{Dažnos klaidos.}
Mastelyje negalima tiesiog lyginti $cm$ su $m$ nepakeitus vienetu. Vektorius nėra tik ilgis: du vienodo ilgio, bet skirtingu krypciu vektoriai nėra lygus.

\subsection{Iliustruojantys pavyzdžiai}
\textbf{1 pavyzdys.} Kambario ilgis plane $6$ cm, mastelis $1:50$. Raskime tikraji ilgį.
\[
6\cdot50=300\ \text{cm}=3\ \text{m}.
\]
\textbf{Atsakymas: $3$ m.}

\textbf{2 pavyzdys.} Tikrasis atstumas $2{,}4$ km. Koks jis žemėlapyje masteliu $1:30000$?
\[
2{,}4\ \text{km}=240000\ \text{cm}.
\]
\[
\frac{240000}{30000}=8\ \text{cm}.
\]
\textbf{Atsakymas: $8$ cm.}

\textbf{3 pavyzdys.} Raskime vektoriaus $\overrightarrow{AB}$ koordinates, kai $A(-2;5)$, $B(4;1)$.
\[
\overrightarrow{AB}=(4-(-2);\ 1-5)=(6;-4).
\]
\textbf{Atsakymas: $(6;-4)$.}

\textbf{4 pavyzdys.} Taska $P(3;-1)$ pastumiame vektoriumi $(-5;2)$.
\[
P'=(3-5;\ -1+2)=(-2;1).
\]
\textbf{Atsakymas: $P'(-2;1)$.}

\paragraph{Pasitikrinkite.}
\begin{enumerate}[label=\alph*)]
\item Mastelis $1:200$. Kiek tikroveje atitinka $7$ cm plane?
\item Raskite $\overrightarrow{CD}$, kai $C(1;1)$, $D(-3;6)$.
\item Taska $(0;4)$ pastumkite vektoriumi $(2;-7)$.
\end{enumerate}

\subsection{Ryšys su kitomis temomis}
Mastelis remiasi proporcijomis \ref{sec:procentai-proporcijos-ir-mastelis}; vektoriai ir transformacijos siejasi su koordinačių plokštuma \ref{sec:transformacijos-ir-simetrijos}, Pitagoro teorema \ref{sec:pitos-teorema} ir erdviniu kunu projekcijomis \ref{sec:erdviniai-kunai-ir-ju-pjuviai}.

\section{Diagramų analizė ir sklaida}\label{sec:diagramu-analize-ir-sklaida}

\subsection{Apibrėžtys}
\begin{apibreztis}
Histograma yra diagrama, kurioje intervalais sugrupuoti duomenys vaizduojami stulpeliais. Kai intervalu ilgiai vienodi, stulpelio aukstis gali rodyti dazni.
\end{apibreztis}
\begin{apibreztis}
Empirinis tankis apibudina, kaip dažniai pasiskirsto intervaluose; kai intervalu ilgiai nevienodi, stulpelio aukstis siejamas su dazniu, padalytu iš intervalo ilgio.
\end{apibreztis}
\begin{apibreztis}
Sukauptasis daznis iki pasirinktos ribos rodo, kiek duomenų reikšmių yra ne didesnes už ta riba. Sukauptasis santykinis daznis gaunamas sukauptaji dazni padalijus iš visu stebejimu skaičiaus.
\end{apibreztis}
\begin{apibreztis}
Isskirtis yra reikšmė, kuri labai skiriasi nuo daugumos duomenų ir gali stipriai paveikti vidurki bei sklaida.
\end{apibreztis}

\subsection{Teoremos ir savybės}
\begin{savybe}
Grupuojant duomenis i intervalus, prarandama dalis tikslumo, bet lengviau pamatyti bendro pasiskirstymo forma.
\end{savybe}
\begin{savybe}
Sukauptieji dažniai niekada nemazeja: judant per intervalus iš kairės i dešinę jie tik dideja arba lieka tokie patys.
\end{savybe}
\begin{savybe}
Staciakampe diagrama su usais leidzia greitai palyginti duomenų centra, sklaida ir išskirtis.
\end{savybe}
\begin{pastaba}
Diagrama yra argumentas tik tada, kai zinome, kaip duomenys surinkti, kokia imtis, kokie intervalai ir kokia skale.
\end{pastaba}

\subsection{Taisyklės ir algoritmai}
\begin{enumerate}[label=\arabic*.]
\item Histogramai sudaryti pasirinkite vienodo ilgio intervalus, suskaiciuokite, kiek reikšmių patenka i kiekviena intervala, ir nubraizykite stulpelius.
\item Sukauptaji dazni skaiciuokite vis pridedami einamojo intervalo dazni prie ankstesniu dazniu sumos.
\item Stačiakampei diagramai raskite penkiu skaičių santrauka: minimuma, $Q_1$, mediana, $Q_3$, maksimuma.
\item Analizuodami diagrama klauskite: kur centras, kokia sklaida, ar yra isskirciu, ar diagrama neiskraipo mastelio.
\end{enumerate}
\paragraph{Dažnos klaidos.}
Histogramoje stulpeliai paprastai lieciasi, nes intervalai eina vienas po kito. Stulpeline diagrama kategorijoms ir histograma intervaliniams duomenims nėra tas pats. Dar viena klaida -- daryti išvadą apie priezasti vien iš diagramos, kuri rodo tik duomenų pasiskirstyma.

\subsection{Iliustruojantys pavyzdžiai}
\textbf{1 pavyzdys.} Duomenys apie mokiniu skaitymo laika minutemis: $12,18,21,25,27,31,34,36,44,48$. Sugrupuokime i intervalus $[10;20)$, $[20;30)$, $[30;40)$, $[40;50)$.
Dazniai:
\[
2,\quad 3,\quad 3,\quad 2.
\]
Sukauptieji dažniai:
\[
2,\quad 5,\quad 8,\quad 10.
\]

\textbf{2 pavyzdys.} Raskime sukauptaji santykini dazni iki $30$ minuciu imtinai, jei iki sio intervalo sukauptasis daznis $5$, o iš viso $10$ stebejimu.
\[
\frac5{10}=0{,}5=50\%.
\]
\textbf{Atsakymas: $50\%$.}

\textbf{3 pavyzdys.} Duomenų penkiu skaičių santrauka yra $4$, $7$, $9$, $13$, $22$. Ką reiškia šie skaičiai?
Minimumas $4$, pirmasis kvartilis $7$, mediana $9$, treciasis kvartilis $13$, maksimumas $22$. Puse duomenų yra tarp $7$ ir $13$.

\textbf{4 pavyzdys.} Duomenys $5,6,6,7,8,9,40$. Kodel $40$ svarbus analizei?
Reiksme $40$ labai nutolusi nuo kitu. Jį padidins vidurki:
\[
\overline{x}=\frac{5+6+6+7+8+9+40}{7}=\frac{81}{7}\approx11{,}57,
\]
nors dauguma reikšmių yra tarp $5$ ir $9$. Mediana yra $7$.

\paragraph{Pasitikrinkite.}
\begin{enumerate}[label=\alph*)]
\item Duomenų dažniai intervaluose yra $4,6,5,3$. Raskite sukauptuosius dažnius.
\item Jei iš viso $30$ stebejimu, o sukauptasis daznis iki ribos yra $18$, koks sukauptasis santykinis daznis?
\item Paaiskinkite, kuo histograma skiriasi nuo skritulines diagramos.
\end{enumerate}

\subsection{Ryšys su kitomis temomis}
Diagramų analizė pratęsia statistikos temą \ref{sec:statistika-imtys-ir-sklaida}. Jį remiasi procentais \ref{sec:procentai-proporcijos-ir-mastelis}, imties samprata, duomenų interpretavimu ir ruosia gilesne statistiniu rodikliu bei koreliacijos temą \ref{sec:statistiniai-rodikliai-ir-koreliacija}.
